\documentclass[11pt]{article}

\usepackage[T1]{fontenc}
\usepackage{amsmath,amssymb}
\usepackage{booktabs}
\usepackage{hyperref}

\title{Classical Particle Mechanics as a Suppes Predicate}
\author{Formalization notes}
\date{\today}

\begin{document}
\maketitle

\begin{abstract}
This note expresses the axiomatic system for classical particle mechanics
of McKinsey, Sugar, and Suppes as a set-theoretic predicate, and gives a
Lean 4 encoding of its model structure. It distinguishes the six axioms
from the additional results discussed by the original authors. The Lean
file formalizes the axiom schema over an explicit abstract signature; it
does not implement real analysis or prove the original paper's embedding
theorems.
\end{abstract}

\section{The Suppes-style predicate}

Fix an ambient signature consisting of a real time domain $T$, a
finite-dimensional real vector space $V$, a force-index set $I$, and the
usual notions of interval, positive real, twice differentiable trajectory,
absolute convergence, and vector sum. A candidate system is the tuple
\[
  \mathcal{M} = (P,T,m,s,f),
\]
where $P$ is the particle set, $m:P\to\mathbb{R}$ is mass,
$s:P\times T\to V$ is position, and
$f:P\times T\times I\to V$ gives the individual forces. Its predicate is
\[
\operatorname{CM}(\mathcal{M}) \quad\Longleftrightarrow\quad
  P_1\land P_2\land P_3\land P_4\land P_5\land P_6,
\]
with the clauses in Table~\ref{tab:axioms}. This is a definition of which
structures count as models, not a claim that every arbitrary tuple
satisfies the laws.

\begin{table}[ht]
\centering
\begin{tabular}{@{}clp{0.65\linewidth}@{}}
\toprule
 & Axiom & Condition \\
\midrule
$P_1$ & Particles & $P$ is finite and nonempty.\\
$P_2$ & Time & $T$ is an interval in $\mathbb{R}$.\\
$P_3$ & Kinematics &
For each $p\in P$, $t\mapsto s(p,t)$ is twice differentiable on $T$.\\
$P_4$ & Mass & For each $p\in P$, $m(p)>0$.\\
$P_5$ & Forces &
For each $(p,t)\in P\times T$, the series
$\sum_{i\in I}f(p,t,i)$ is absolutely convergent.\\
$P_6$ & Dynamics &
For each $(p,t)\in P\times T$,
\[
  \sum_{i\in I} f(p,t,i)
  = m(p)\,\frac{\partial^2 s}{\partial t^2}(p,t).
\]\\
\bottomrule
\end{tabular}
\caption{The axioms defining an $n$-dimensional system of particle
mechanics. The vector space $V$ is $\mathbb{R}^n$.}
\label{tab:axioms}
\end{table}

\section{Reading the axioms and the reported results}

The kinematic and dynamical clauses work together: $P_3$ supplies
acceleration as the second time derivative of position, $P_5$ makes the
total force well-defined independently of the order of summation, and
$P_6$ equates that force with mass times acceleration. The original paper
does not assume Newton's first or third law as axioms. It reports the
first law as a consequence of its system and proves an embedding result
relating its models to Newtonian systems; those results are not formalized
in the Lean file here \cite{mckinsey1953axiomatic}.

In the Suppes approach, a set-theoretic predicate is a formula on
structures that selects the models of a theory. Da Costa and Krause
emphasize that transportability is relevant when the intent is to capture
all structures of a kind, but need not be imposed when a predicate is
intended to select only a particular class \cite{daCostaKrause2020}.
Accordingly, $\operatorname{CM}$ is stated relative to the ambient
signature and to the chosen class of particle systems; no claim of
transportability is made here. Suppes's broader account of scientific
structures is a methodological reference for this representation
\cite{suppes2002representation}.

\section{Lean encoding and scope}

The companion file \texttt{\detokenize{_lean/ClassicalMechanics.lean}}
defines a signature carrying the semantic operations and predicates for
time intervals, positive masses, differentiability, summability, and force
summation. A \texttt{System} supplies particles and the functions $m$,
$s$, acceleration, and $f$, together with witnesses of $P_1$--$P_6$.
The theorem \texttt{finiteNonemptyParticles} packages $P_1$, and
\texttt{secondLawHolds} exposes $P_6$ as a checked Lean theorem.

This is a typed, machine-checked formalization of the structural axiom
schema, not a development of calculus over $\mathbb{R}$. The interpretation
of the signature's predicates and operations must be supplied by a
mathematical library or a concrete model before numerical or analytic
consequences can be proved. The result is therefore deliberately narrower
than a full formalization of the cited paper.

\bibliographystyle{plain}
\bibliography{references}

\end{document}
